\documentclass[sigconf]{acmart}
\usepackage{makecell}
\AtBeginDocument{%
  \providecommand\BibTeX{{%
    Bib\TeX}}}

\setcopyright{acmlicensed}
\copyrightyear{2018}
\acmYear{2018}
\acmDOI{XXXXXXX.XXXXXXX}

\acmConference[Conference acronym 'XX]{Make sure to enter the correct
  conference title from your rights confirmation email}{June 03--05,
  2018}{Woodstock, NY}
\acmISBN{978-1-4503-XXXX-X/2018/06}


\usepackage{multirow} 
\newcommand{\HU}[1]{{\color{cyan} [{#1}]}} 
\newcommand{\UC}[1]{{\color{pink} [{#1}]}} 
\newcommand{\NR}[1]{{\color{red} [{#1}]}} 
\newcommand{\Changed}[1]{{\color{blue} [{#1}]}}

\begin{document}


\title{Do We Really Need Multimodal Emotion Language Models Larger Than 1B Parameters?}

\settopmatter{authorsperrow=3}

\author{Kaiwen Zheng}
\authornote{Equal contribution.}
\affiliation{%
  \institution{University of Glasgow}
  \country{United Kingdom}
}
\email{k.zheng.1@research.gla.ac.uk}

\author{Junchen Fu}
\authornotemark[1]
\affiliation{%
  \institution{University of Glasgow}
  \country{United Kingdom}
}
\email{j.fu.3@research.gla.ac.uk}

\author{Wenhao Deng}
\affiliation{%
  \institution{University of Glasgow}
  \country{United Kingdom}
}
\email{w.deng.1@research.gla.ac.uk}

\author{Hu Han}
\affiliation{%
  \institution{Institute of Computing Technology, Chinese Academy of Sciences}
  \country{China}
}
\email{hanhu@ict.ac.cn}

\author{Joemon M. Jose}
\affiliation{%
  \institution{University of Glasgow}
  \country{United Kingdom}
}
\email{Joemon.Jose@glasgow.ac.uk}

\author{Xuri Ge}
\authornote{Corresponding Author. Xuri Ge is also a member of the Qingdao Key Laboratory of Trustworthy Artificial Intelligence in Qingdao, China.}
\affiliation{%
  \institution{School of Artificial Intelligence, }
  \institution{Shandong University}
  \country{China}
}
\email{xuri.ge@sdu.edu.cn}

\renewcommand{\shortauthors}{Zheng et al.}

\begin{abstract}
Recent advances in multimodal large language models (MLLMs) have significantly improved the performance of multimodal emotion recognition (MER) and enabled interpretable description generation by jointly modeling video, audio, and language, etc. However, these performance improvements are often accompanied by an increase in model parameter size (e.g, at least 7B), which simultaneously incurs high computational costs and reduces inference efficiency, thereby hindering real-time deployment on resource-constrained platforms such as robots and mobile devices. This raises a fundamental question: do we really need the multimodal MER model larger than 1B parameters for high-quality MER? 

In this paper, we challenge the assumption that larger models are inherently necessary and proposes a lightweight MER framework (called \textbf{Light-MER}), which achieves better and faster multimodal sentiment understanding and recognition through knowledge distillation.
It can transfer knowledge from a strong, large-scale teacher model to a lightweight sub-billion-parameter student model, aiming to preserve rich multimodal emotion reasoning and recognition while substantially improving deployment efficiency.
Specifically, we introduce two new optimization strategies to enhance knowledge transfer: (1) a new optimal transport loss that combines Sliced Wasserstein Distance with hidden-state alignment, and (2) a new multi-reward optimization strategy based on GRPO that balances MER performance and efficiency, aimed at further enhancing the learning capabilities of student models. Extensive experiments on nine benchmark datasets demonstrate that Light-MER achieves state-of-the-art performance while significantly improving inference efficiency. This highlights the strong potential of small multimodal emotion language models for future research. Code is available at https://github.com/GAIR-Lab/Light-MER.
\end{abstract}

\begin{CCSXML}
<ccs2012>
   <concept>
       <concept_id>10002951.10003317.10003338.10003341</concept_id>
       <concept_desc>Information systems~Language models</concept_desc>
       <concept_significance>500</concept_significance>
       </concept>
 </ccs2012>
\end{CCSXML}

\ccsdesc[500]{Information systems~Language models}

\keywords{Multimodal Emotion Recognition, Multimodal Large Language Models, Knowledge Distillation, Sliced Wasserstein Distance}

\maketitle
\begin{figure}[t]
    \centering
    \includegraphics[width=\linewidth]{Image/motivation2.pdf}
    \vspace{-2em}
    \caption{Large multimodal emotion models achieve strong performance but are difficult to deploy at the edge. Our goal is to preserve most of the performance of an 8B teacher while moving the deployment model below 1B parameters.}
    \label{fig:motivation}
    \vspace{-1.5em}
\end{figure}
\section{Introduction}

Multimodal Emotion Recognition (MER) serves as a cornerstone for developing perceptive intelligent systems, demonstrating immense potential across domains such as social robotics \cite{MERRob}, healthcare assistance \cite{MERhealthcare}, and smart education\cite{gan2022iot}. Given the inherent complexity of affective expressions, linguistic cues alone often fail to yield robust inferences. In contrast, facial expressions, acoustic signals, scene dynamics, and textual context provide critical, complementary evidence. Consequently, a practical MER framework must facilitate the synergistic modeling of visual, audio, and textual modalities while maintaining computational efficiency to accommodate real-world resource constraints.

Existing MER methods can be broadly divided into discriminative and generative paradigms. Traditional MER \cite{TowardXuri} is mostly discriminative, where multimodal features are mapped to predefined emotion categories through task-specific classifiers. Although effective, these methods are limited to closed-set prediction and offer little interpretability. Recently, the rapid development of multimodal large language models (MLLMs) \cite{Yin_2024} has driven MER toward a more generative paradigm. Instead of predicting only a label, generative MER models \cite{zhao2026generative} can describe emotional states in natural language and provide richer evidence for their predictions. For example, AffectGPT \cite{affectgpt} shows that multimodal generators can capture more nuanced and compositional affective states than conventional classifiers. This shift is particularly important because human emotions are often subtle, ambiguous, and co-occurring.


Despite this progress, as illustrated in Figure \ref{fig:motivation}, current generative MER systems based on pre-trained MLLMs remain heavily reliant on increasingly large model scales—typically requiring at least 7B parameters \cite{EVLM2025, Emotion-LLaMA}—to achieve strong recognition and generation performance.
Although such scaling improves benchmark results, it also introduces substantial computational overhead, memory consumption, and inference latency, which severely limit deployment on robots, smartphones, and other edge devices. More fundamentally, generative MER requires the model to both integrate heterogeneous multimodal evidence and produce coherent emotion-aware descriptions, making efficient deployment substantially more challenging than classification-based MER. This naturally raises a key question: \textit{do we really need multimodal emotion language models larger than 1B parameters for high-quality MER?}

We argue that this is not necessary, provided that the multimodal reasoning abilities of large generative models can be effectively distilled into compact deployment models.
This motivates us to revisit MER from the perspective of efficient multimodal knowledge transfer. As illustrated in Figure~\ref{fig:motivation}, practical MER systems require a better balance among model size, accuracy, interpretability, and inference efficiency than current large-model solutions can provide. A natural way to achieve this goal is to first leverage a strong large-scale MER model as a source of rich multimodal knowledge, and then transfer such knowledge into a lightweight model suitable for deployment. In this way, as shown in Figure \ref{fig:motivation}, the large model serves as a teacher, while the compact model acts as a student.
 
The main challenge lies in how to preserve the teacher’s internal multimodal reasoning process during compression. Existing output-level distillation \cite{gu2026mini} based on KL divergence \cite{KL2014} mainly constrains the final token distribution, but does not preserve the latent reasoning structure that supports emotion understanding. Likewise, simple hidden-state regression losses such as MSE \cite{sun2019} ignore the distributional geometry of the teacher representation space. For MER, where subtle cues from voice, face, motion, and text must be integrated into coherent latent affective semantics, preserving such internal geometry is crucial for effective transfer. Therefore, a high-quality compact MER model requires not only model compression, but also a distillation framework that better preserves multimodal representation structure and generative capability.

To this end, we propose \textbf{Light-MER}, a knowledge distillation framework with multi-reward GRPO refinement for better and faster multimodal emotion understanding and recognition. Light-MER transfers knowledge from a strong large-scale MLLM to a lightweight sub-billion-parameter student model, aiming to retain rich multimodal emotion reasoning while substantially improving deployment efficiency. Specifically, our framework introduces two key optimization strategies. 
First, we define SWD-H as a hidden-state alignment objective that utilizes the Sliced Wasserstein Distance to minimize the distance between teacher and student latent distributions. Unlike point-wise alignment, SWD-H treats hidden representations as empirical distributions, facilitating an optimal transport-based mapping that is robust to high-dimensional manifold shifts.
Second, we further refine the distilled student with a new GRPO-based optimization strategy (called M-GRPO), designed with simultaneous constraints on multi-reward performance and efficiency for generative MER. 
We hypothesize that while hidden-state alignment (SWD-H) captures the teacher’s latent geometry, it does not inherently optimize for generative output quality. Therefore, we introduce M-GRPO as a critical refinement stage to align the model’s generated narratives with multi-faceted rewards, specifically focusing on the nuance of MER.


Our contributions are threefold:
\begin{itemize}
    \item We revisit multimodal emotion recognition from an efficiency perspective and ask whether high-quality generative MER truly requires deployment models larger than 1B parameters.
    \item We propose Light-MER, a compact teacher-student framework that transfers multimodal emotion understanding and recognition ability from a large generative MLLM to a lightweight student model for deployment. 
    \item We introduce SWD-H, a hidden-state alignment objective based on Sliced Wasserstein Distance that explicitly preserves the geometry of multimodal latent structures, thereby improving the effectiveness of knowledge distillation.
    \item We propose M-GRPO, a new GRPO-based optimization strategy with simultaneous constraints on multi-reward performance and efficiency, to further improve generation quality and student learning for generative MER.
\end{itemize}
Extensive experiments on nine benchmark datasets demonstrate that Light-MER not only significantly improves inference efficiency, but also achieves state-of-the-art performance, showing that high-quality generative MER does not necessarily require deployment models larger than 1B parameters.

\section{Related Work}
\noindent \emph{Generative Multimodal Emotion Recognition.}
Early MER is mainly formulated as a discriminative classification problem, where multimodal features are fused and mapped to a predefined label space through task-specific classifiers.
Representative methods include multi-task learning with modality-specific encoders \cite{zheng-etal-2023-facial,fu2024iisan}, graph-based cross-modal reasoning \cite{Li_2023}, and attention-based fusion \cite{shi-huang-2023-multiemo}.
While effective in closed-label settings, these methods are limited in modeling ambiguity, co-occurrence, and fine-grained affective intensity.
Recent surveys summarize progress in datasets, fusion strategies, and model architectures \cite{biomimetics10070418, AV2024102218}.

With the rise of multimodal large language models (MLLMs) \cite{openai2024,yang2025,geminiteam2025,fu2025efficient,fu2026differentiable,fu2024exploring,he2025double,fu2026streamaware}, MER has shifted toward a generative formulation, where models describe emotions in natural language and support open-vocabulary affective reasoning.
Emotion-LLaMA \cite{Emotion-LLaMA} joints video, audio, and text-based emotion understanding with in a unified large language model.
AffectGPT \cite{affectgpt} further combines multimodal encoders with a 7B language decoder and introduces the MER-Caption+ dataset.
A recent survey \cite{shou2025} likewise highlights the growing role of MLLMs in emotion recognition and reasoning.
However, existing generative MER systems typically rely on deployment models of at least 7B parameters to achieve strong performance \cite{affectgpt,Emotion-LLaMA,su2023}, leading to substantial memory and latency costs that hinder practical deployment on resource-constrained devices.
Our work follows the generative MER paradigm, but challenges the assumption that high-quality multimodal emotion understanding must depend on a large deployment model.


\noindent \emph{Knowledge Distillation and Optimal Transport Alignment.}
Knowledge distillation \cite{hinton2015} is a standard way to transfer the capability of large models to compact students.
Most prior methods \cite{gu2026mini,agarwal2024} rely on softened cross-entropy or KL divergence over output distributions, such as DistilBERT \cite{sanh2020}.
However, logit-level distillation is less suitable for generative MER, where the goal is to preserve multimodal reasoning behavior rather than only next-token probabilities.
In our setting, teacher outputs are often highly peaked, which weakens the supervision signal beyond the top token.

Recent work therefore explores hidden-state distillation, including layer-wise alignment in TinyBERT \cite{jiao2020}, reverse-KL distillation in MiniLLM \cite{gu2026mini}, and recent analyses highlight hidden-state alignment as a promising direction \cite{dasgupta2025improving}.
However, pointwise objectives such as MSE or cosine distance do not explicitly preserve representation geometry, while CKA-based alignment \cite{zhou2024} only partially addresses this issue.
This limitation is especially important for multimodal generation, where hidden states organize evidence across modalities before decoding.

Optimal transport (OT) offers a principled approach to geometry-aware alignment \cite{cui2024,lv2024,zhuang2025frequency}. Compared with KL-based objectives, Wasserstein distances are more robust under partial distribution mismatch. However, existing OT variants such as Sinkhorn OT \cite{NIPS2013_af21d0c9}, position-aware OT \cite{PAOT}, and unbalanced OT \cite{UnbalancedOT} either require iterative optimization or introduce additional coupling and balancing terms, increasing computational cost in repeated hidden-state matching. We therefore adopt Sliced Wasserstein Distance (SWD) \cite{kolouri2019}, which enables efficient OT approximation via random one-dimensional projections and sorting. SWD retains the geometric benefits of OT with much lower overhead and has shown effectiveness in recent compression settings \cite{quetu2025}. Accordingly, we use SWD to align teacher and student hidden-state distributions for generative MER.
Further details on alternative OT variants are provided in the supplementary material.

\noindent \emph{Reinforcement Learning for Post-distillation Refinement.}
Reinforcement learning is widely used to improve generation quality in large language models.
RLHF \cite{ouyang2022} and PPO-based optimization \cite{schulman2017} are effective, but require an additional value model or critic.
GRPO \cite{shao2024} removes the critic by estimating advantages from multiple responses to the same prompt and normalizing rewards within each group.
This substantially simplifies training and has been adopted in recent reasoning-oriented models such as DeepSeekMath \cite{shao2024} and DeepSeek-R1 \cite{Guo_2025}.

We adapt GRPO as a post-distillation refinement stage for generative MER. Instead of a single reward, we use a multi-reward objective to improve emotion accuracy, conciseness, information density, completeness, and low repetition. In this way, reinforcement learning complements representation-level distillation by improving output quality and efficiency.

\begin{figure*}[t!]
    \centering
    % \vspace{-1.3em}
    \includegraphics[width=\textwidth]{Image/model3.pdf}
    \vspace{-2em}
    \caption{An overview of our model. A frozen 8B teacher supervises a sub-1B student through SWD hidden-state alignment (SWD-H), and the distilled student is further refined with multi-reward GRPO for concise emotion description generation.}
    \label{fig:model}
    \vspace{-1em}
\end{figure*}

\section{Method}

\subsection{Problem Formulation}

Given a multimodal sample $\mathcal{X}=\{x^{\mathrm{frm}}/x^{\mathrm{face}},x^{\mathrm{aud}},x^{\mathrm{sub}}\}$ and an instruction $q$, the model generates an emotion description $\mathcal{Y}=\{y_1,\ldots,y_L\}$. The visual stream supports either full video frames $x^{\mathrm{frm}}$ or face-cropped clips $x^{\mathrm{face}}$; as facial regions carry the most salient affective cues, we adopt face crops in all experiments. We consider a teacher model $\mathcal{T}$ and a student model $\mathcal{S}$. The teacher is a high-capacity multimodal emotion language model used only during training, while the student is the deployment model. Our goal is to train $\mathcal{S}$ to (i) preserve the teacher's multimodal emotion understanding ability through hidden-state distillation, and (ii) produce shorter and cleaner descriptions through policy optimization.
The standard autoregressive objective is
\begin{equation}
\mathcal{L}_{\mathrm{CE}}
=
-\sum_{l=1}^{L}\log p_{\theta}(y_l\mid x,q,y_{[1:l-1]}),
\end{equation}
where $\theta$ denotes the student parameters.

Building a distillation framework on top of this objective requires two design choices: (i) at which representation level should teacher and student be aligned (output logits, intermediate hidden states, or cross-layer relations), and (ii) which divergence measure should drive the alignment. We address (i) through an empirical analysis of hidden-state and output distributions in Section~\ref{subsec:observation}, and (ii) through the SWD formulation in Section~\ref{subsec:swd}. The student is then further refined with GRPO-based multi-reward policy optimization in Section~\ref{M-GRPO}.

\subsection{Model Architecture}

Figure~\ref{fig:model} summarizes the overall framework of our proposed Light-MER.
The teacher consists of Qwen3-8B \cite{yang2025} as the language decoder,
CLIP-ViT-Large-Patch14 \cite{radford2021} as the visual encoder, and
HuBERT-Large \cite{hsu2021hubert} as the acoustic encoder.
It is pre-trained on the MER-Caption+ \cite{affectgpt} corpus and then frozen
during student distillation.
The student replaces the teacher with a deployment-scale architecture:
Qwen3-0.6B \cite{yang2025} as the language decoder,
CLIP-ViT-Base-Patch16 \cite{radford2021} as the visual encoder, and
HuBERT-Base \cite{hsu2021hubert} as the acoustic encoder.
The LLM branch is updated through LoRA while the multimodal projectors
and fusion layers are trained jointly with the student decoder.
The central design principle is that the teacher should remain large
enough to provide rich multimodal hidden-state supervision, while the
student should remain small enough for practical deployment.

\paragraph{Multimodal Fusion.}
Each modality stream $m\!\in\!\{\mathrm{face},\mathrm{aud}\}$ is first
encoded by a corresponding frozen encoder into a temporal sequence
$\mathbf{Z}^{(m)}\!=\!\{\mathbf{z}_1^{(m)},\ldots,\mathbf{z}_{T_m}^{(m)}\}$,
then compressed into a single vector via learned temporal pooling
$\bar{\mathbf{z}}^{(m)}\!=\!\sum_t a_t^{(m)}\mathbf{z}_t^{(m)}$
and projected into the LLM input space by a linear mapping
$\mathbf{h}^{(m)}\!=\!\mathbf{W}_m\bar{\mathbf{z}}^{(m)}$.
The two streams are merged through a learnable MLP-based gating mechanism:
\begin{equation}
\mathbf{h}^{(\mathrm{mm})}
=
\beta_{\mathrm{v}}\,\mathbf{h}^{(\mathrm{face})}
+
\beta_{\mathrm{aud}}\,\mathbf{h}^{(\mathrm{aud})}, 
\end{equation}
\begin{equation}
[\beta_{\mathrm{face}},\beta_{\mathrm{aud}}]
\!=\textbf{MLP}[\mathbf{h}^{(\mathrm{face})}: \mathbf{h}^{(\mathrm{aud})}],
\end{equation}
where [:] indicates the concatenation.
The LLM receives the fused multimodal token $\mathbf{h}^{(\mathrm{mm})}$
together with separate face tokens, audio tokens, the text subtitle,
and the instruction prompt, and decodes via the Qwen3-0.6B student
with LoRA adaptation.
All reported results use this gating-based fusion with face inputs.

\subsection{Empirical Observation: Hidden States vs.\ Output Distributions}
\label{subsec:observation}

Before distillation begins, we compare the 8B and 0.6B models that have been independently trained on the MER task (Figure~\ref{fig:hidden_logit_analysis}). We aggregate statistics across subsets sampled from the nine-benchmark suite and examine output probabilities and last-layer hidden states.

Output probabilities are extremely peaked for both models (Figure~\ref{fig:hidden_logit_analysis}a). The 8B model places 0.980 of its probability on the top-1 token. Tokens ranked 2--20 collectively carry only 0.016 of the 8B model's mass. When the teacher's distribution is this peaked, KL-based distillation behaves similarly to hard-label cross-entropy training (the gradient of KL divergence at each token is weighted by the teacher's probability, so the top-1 token dominates), leaving little room for the student to learn from the teacher's output distribution beyond the top prediction.

Across the same aggregated data, hidden-state activations exhibit multi-modal structure in 1016 of 1024 dimensions, or 99.2\% (Figure~\ref{fig:hidden_logit_analysis}c). A representative sample (Figure~\ref{fig:hidden_logit_analysis}b) illustrates the pattern: the 8B and 0.6B models place modes at different positions and heights. Despite producing nearly identical outputs, the two models organize their internal representations differently. OT-based objectives can account for the distance between these modes, while pointwise losses such as MSE only penalize per-position differences without considering the overall distribution shape. The SWD-H formulation below builds on this observation. A visual comparison before and after SWD-H training is shown in Figure~\ref{fig:hidden_logit_analysis}(b,\,d) and discussed in Section~\ref{sec:ablation}.

\begin{figure*}[t]
    \centering
    \includegraphics[width=0.9\textwidth]{Image/figure3_ot_motivation.pdf}
    \vspace{-1.0em}
    \caption{Hidden-state and output-probability analysis aggregated across subsets evaluation suites.
    (a) Top-20 token probabilities.
    (b) A representative hidden-state dimension before distillation. Teacher and student are both multi-modal with different mode locations.
    (c) Mode-count histogram across all 1024 dimensions. 99.2\% of dimensions contain two or more modes.
    (d) The same dimension as (b) after SWD-H distillation. Teacher and student modes are more aligned.}
    \label{fig:hidden_logit_analysis}
    \vspace{-0.6em}
\end{figure*}

\subsection{Sliced Wasserstein Distillation on Hidden States (SWD-H)}
\label{subsec:swd}

Our distillation target is the last-layer hidden state of the LLM rather than the output logits. Let $\mathbf{H}^{T}\in\mathbb{R}^{L\times4096}$ and $\mathbf{H}^{S}\in\mathbb{R}^{L\times1024}$ denote the teacher and student hidden states at the supervised decoding positions. To compare them in the same space, we use an SWD projector composed of a frozen orthogonally initialized teacher projection, an identity-initialized student projection, and LayerNorm:
\begin{equation}
\tilde{\mathbf{H}}^{T}=\mathrm{LN}(\mathbf{H}^{T}\mathbf{W}_{T}),
\qquad
\tilde{\mathbf{H}}^{S}=\mathrm{LN}(\mathbf{H}^{S}\mathbf{W}_{S}),
\end{equation}
where $\mathbf{W}_T\in\mathbb{R}^{4096\times1024}$ is frozen and $\mathbf{W}_S\in\mathbb{R}^{1024\times1024}$ is identity-initialized. Following standard autoregressive training, the cross-entropy loss is computed only at positions whose next token belongs to the answer sequence. We apply the same shifted mask as cross-entropy to the hidden-state distillation: for each position $i$ in the sequence, SWD alignment is active if and only if $y_{i+1} \neq -100$ (Ignore Index), where $y$ denotes the label sequence.

We then sample $R=100$ random unit directions $\{\boldsymbol{\theta}_r\}_{r=1}^{R}$ in $\mathbb{R}^{1024}$ and compute one-dimensional projections for teacher and student. The SWD-H loss is
\begin{equation}
\mathcal{L}_{\mathrm{SWD}}
=
\frac{1}{R}
\sum_{r=1}^{R}
\left\|
\mathrm{sort}\!\left(\tilde{\mathbf{H}}^{T}_{M}\boldsymbol{\theta}_r\right)
-
\mathrm{sort}\!\left(\tilde{\mathbf{H}}^{S}_{M}\boldsymbol{\theta}_r\right)
\right\|_2^2,
\end{equation}
where $M$ denotes the answer-token mask. In one dimension, the optimal transport plan between two empirical distributions of equal size $n$ is the quantile coupling: sorting both sets of values and matching them index-by-index is provably optimal \cite{peyre2020}. The $p$-Wasserstein distance therefore reduces to
\begin{equation}
W_p^p(\alpha,\beta)
=
\frac{1}{n}\sum_{i=1}^{n}\bigl|\alpha_{(i)}-\beta_{(i)}\bigr|^{p},
\end{equation}
where $\alpha_{(1)}\!\leq\!\cdots\!\leq\!\alpha_{(n)}$ and $\beta_{(1)}\!\leq\!\cdots\!\leq\!\beta_{(n)}$ are the sorted projections ($p{=}2$ in our case). This closed-form solution gives $\mathcal{L}_{\mathrm{SWD}}$ an $O(Rn\log n)$ complexity dominated by sorting, compared with $O(Kn^{2})$ for $K$-iteration Sinkhorn, and eliminates the need for entropic regularization $\epsilon$.

The total distillation objective is
\begin{equation}
\mathcal{L}
=
\mathcal{L}_{\mathrm{CE}}
+
\alpha_t\,\mathcal{L}_{\mathrm{SWD}},
\qquad
\alpha_t
=
\lambda_{\mathrm{SWD}}
\min\!\left(\frac{t}{5000},1\right),
\end{equation}
where $\lambda_{\mathrm{SWD}}=1.0$ and the SWD-H weight $\alpha_t$ is linearly warmed up over the first 5000 steps. Across ten distillation variants evaluated in this study, SWD-H is the best overall strategy.



\subsection{Multi-Reward GRPO-based (M-GRPO) Description Refinement} \label{M-GRPO}

Distillation preserves the teacher's emotion reasoning ability, but it also transfers the teacher's verbose generation pattern, a common issue in large-scale affective models that produce over-elaborate emotional descriptions.
In our setting, the model output consists of a fixed prefix that restates the subtitle (e.g., ``In the text, the caption reads: \ldots'') followed by a free-form emotion analysis.
We apply GRPO only to the emotion analysis segment; the fixed template prefix is excluded from the policy optimization so that the model is not rewarded for simply copying it. Since emotion F1 alone cannot control verbosity, informativeness, sentence completeness, or repetition, we use a composite reward with complementary components.
We define the total reward as
\begin{equation}
r
=
3.0\,r_{\mathrm{F1}}
+
3.0\,r_{\mathrm{len}}
+
2.0\,r_{\mathrm{dens}}
+
2.0\,r_{\mathrm{comp}}
+
0.5\,r_{\mathrm{rep}},
\end{equation}
where $r_{\mathrm{F1}}$ measures average emotion F1 across five emotion wheels against the ground truth, $r_{\mathrm{len}}$ encourages a compact emotion analysis and decays for longer outputs, $r_{\mathrm{dens}}$ rewards emotion-keyword density, $r_{\mathrm{comp}}$ rewards completed sentences ending with punctuation, and $r_{\mathrm{rep}}$ penalizes repetition through the unique-word ratio.


\begin{table*}[t]
    \centering
    \small
    \setlength{\tabcolsep}{4pt}
    \renewcommand{\arraystretch}{0.9}
    \caption{Main results on the evaluation suite. This table follows the layout of AffectGPT Table~2. Primary metrics are reported for each dataset: HIT for basic emotion, WAF for sentiment, and $F_s$ for OV-MERD+. AffectGPT baseline numbers are reproduced using the official open-source code.}
    \vspace{-1em}
    \label{tab:main_results}
    \resizebox{\textwidth}{!}{
    \begin{tabular}{lr|ccc|cccc|cccc|c|c}
        \toprule
        \multirow{2}{*}{Model} & \multirow{2}{*}{LLM Param} & \multicolumn{3}{c|}{Modality} & \multicolumn{4}{c|}{Basic Emotion} & \multicolumn{4}{c|}{Sentiment} & \multicolumn{1}{c|}{Fine-grained} & \multirow{2}{*}{Mean} \\
        & & A & V & T & MER2023 & MER2024 & MELD & IEMOCAP & MOSI & MOSEI & SIMS & SIMS v2 & OV-MERD+ & \\
        \midrule
        \multicolumn{15}{l}{\textbf{Audio+Text MLLMs}} \\
        OneLLM \cite{OneLLM}       & 7B & $\checkmark$ & $\times$     & $\checkmark$ & 25.52 & 17.21 & 28.32 & 33.44 & 64.01 & 54.09 & 63.39 & 61.98 & 22.25 & 41.14 \\
        SECap \cite{SECap}        & 7B & $\checkmark$ & $\times$     & $\checkmark$ & 40.95 & 52.46 & 25.56 & 36.92 & 55.76 & 54.18 & 59.51 & 57.41 & 36.97 & 46.64 \\
        PandaGPT \cite{su2023}     & 7B & $\checkmark$ & $\times$     & $\checkmark$ & 33.57 & 39.04 & 31.91 & 36.55 & 66.06 & 61.33 & 62.93 & 58.88 & 31.33 & 46.84 \\
        Qwen-Audio \cite{Qwen-Audio}  & 7B & $\checkmark$ & $\times$     & $\checkmark$ & 41.85 & 31.61 & 49.09 & 35.47 & 70.09 & 46.90 & 70.73 & 65.26 & 32.36 & 49.26 \\
        SALMONN \cite{SALMONN}     & 13B & $\checkmark$ & $\times$     & $\checkmark$ & 55.53 & 45.38 & 45.62 & \underline{46.84} & \underline{81.00} & 67.03 & 68.69 & 65.93 & 45.00 & 57.89 \\
        AffectGPT(Qwen3-8B) & 8B & $\checkmark$ & $\times$ & $\checkmark$ & \underline{59.28} & \underline{56.31} & \textbf{54.03} & 46.69 & \textbf{81.73} & \underline{78.56} & \underline{77.77} & \underline{78.85} & \underline{52.39} & \underline{65.07} \\
        \textbf{Light-MER(Ours)}  & \textbf{0.6B} & $\checkmark$ & $\times$ & $\checkmark$ & \textbf{59.33} & \textbf{58.30} & \underline{53.10} & \textbf{52.18} & 80.48 & \textbf{79.95} & \textbf{80.74} & \textbf{79.80} & \textbf{53.75} & \textbf{66.40} \\
        \midrule
        \multicolumn{15}{l}{\textbf{Video+Text MLLMs}} \\
        Otter \cite{Otter}           & 9B & $\times$     & $\checkmark$ & $\checkmark$ & 16.41 & 14.65 & 22.57 & 29.08 & 52.89 & 50.44 & 57.56 & 53.12 & 16.63 & 34.82 \\
        Video-LLaVA  \cite{video-llava}    & 7B & $\times$     & $\checkmark$ & $\checkmark$ & 36.93 & 30.25 & 30.73 & 38.95 & 56.37 & 61.64 & 53.28 & 57.45 & 34.00 & 44.40 \\
        PandaGPT \cite{su2023}        & 7B & $\times$     & $\checkmark$ & $\checkmark$ & 39.13 & 47.16 & 38.33 & 47.21 & 58.50 & 64.25 & 62.07 & 65.25 & 35.07 & 50.77 \\
        Video-ChatGPT \cite{video-chatgpt}   & 7B & $\times$     & $\checkmark$ & $\checkmark$ & 44.86 & 46.80 & 37.33 & \textbf{56.83} & 54.42 & 63.12 & 64.82 & 65.80 & 39.80 & 52.64 \\
        % VideoChat2       & 7B & $\times$     & $\checkmark$ & $\checkmark$ & 33.67 & 54.50 & 36.64 & 48.70 & 66.84 & 54.32 & 69.49 & 70.66 & 39.21 & 52.67 \\
        LLaMA-VID  \cite{llamavid}      & 7B & $\times$     & $\checkmark$ & $\checkmark$ & 50.72 & 57.60 & 42.75 & 46.02 & 61.78 & 63.89 & 69.35 & 67.48 & 45.01 & 56.07 \\
        VideoChat  \cite{videochat}      & 7B & $\times$     & $\checkmark$ & $\checkmark$ & 48.73 & 57.30 & 41.11 & 48.38 & 65.13 & 63.61 & 69.52 & 72.14 & 44.52 & 56.71 \\
        Chat-UniVi \cite{chatunivi}      & 7B & $\times$     & $\checkmark$ & $\checkmark$ & 57.62 & \textbf{65.67} & 45.61 & 52.37 & 54.53 & 63.18 & 68.15 & 66.36 & 48.00 & 57.94 \\
        mPLUG-Owl  \cite{mPLUG}      & 7B & $\times$     & $\checkmark$ & $\checkmark$ & 56.86 & \underline{59.89} & 49.11 & \underline{55.54} & 72.40 & 72.91 & 72.13 & 75.00 & 48.18 & 62.45 \\
        AffectGPT(Qwen3-8B) & 8B & $\times$ & $\checkmark$ & $\checkmark$ & \underline{58.70} & 55.28 & \underline{50.51} & 48.59 & \textbf{83.11} & \underline{75.50} & \underline{78.15} & \underline{80.49} & \underline{49.59} & \underline{64.44}\\
        \textbf{Light-MER(Ours)}  & \textbf{0.6B} & $\times$ & $\checkmark$ & $\checkmark$ & 58.77 & 56.31 & \textbf{51.51} & 47.38 & \underline{82.39} & \textbf{79.36} & \textbf{82.80} & \textbf{81.32} & \textbf{51.51} & \textbf{65.71} \\
        \midrule
        \multicolumn{15}{l}{\textbf{Audio+Video+Text MLLMs}} \\
        PandaGPT  \cite{su2023}       & 7B & $\checkmark$ & $\checkmark$ & $\checkmark$ & 40.21 & 51.89 & 37.88 & 44.04 & 61.92 & 67.61 & 68.38 & 67.23 & 37.12 & 52.92 \\
        Emotion-LLaMA \cite{Emotion-LLaMA}   & 7B & $\checkmark$ & $\checkmark$ & $\checkmark$ & 59.38 & 73.62 & 46.76 & 55.47 & 66.13 & 67.66 & 78.32 & 77.23 & 52.97 & 64.17 \\
        AffectGPT \cite{affectgpt} & 7B & $\checkmark$ & $\checkmark$ & $\checkmark$ & \textbf 69.74 & 67.08 & \underline{56.05} & 53.51 & \textbf{83.49} & \underline{78.06} & 82.79 & 80.15 & 57.03 & 69.77 \\
        AffectGPT(Qwen3-8B)& 8B & $\checkmark$ & $\checkmark$ & $\checkmark$ & \textbf{76.72} & \underline{79.68} & 55.99 & \underline{59.32} & 78.46 & \textbf{78.48} & \underline{88.30} & \underline{86.94} & \underline{61.47} & \underline{73.93} \\
        \textbf{Light-MER(Ours)}  & \textbf{0.6B} & $\checkmark$ & $\checkmark$ & $\checkmark$ & \underline{75.66} & \textbf{80.04} & \textbf{57.02} & \textbf{61.46} & \underline{80.04} & 77.23 & \textbf{89.05} & \textbf{87.74} & \textbf{63.28} & \textbf{74.61} \\
        \bottomrule
    \end{tabular}}
    \vspace{-0.5em}
\end{table*}

\begin{table*}[t]
    \centering
    \tiny
    \caption{Full-pipeline inference profiling.}
    \vspace{-1em}
    \fontsize{8.5}{8.5}\selectfont
    \renewcommand\tabcolsep{5pt}  
    \label{tab:inference-profile}
    \begin{tabular}{l c c c c cc cc}
    \toprule
    & & & & & \multicolumn{2}{c}{Direct} & \multicolumn{2}{c}{Descriptive} \\
    \cmidrule(lr){6-7} \cmidrule(lr){8-9}
    Model
      & \makecell{Total\\Params}
      & \makecell{FLOPs\\(G)}
      & \makecell{Peak\\Mem.}
      & Compress.
      & \makecell{Time\\(s/samp.)}
      & \makecell{Words\\/samp.}
      & \makecell{Time\\(s/samp.)}
      & \makecell{Words\\/samp.} \\
    \midrule
    Teacher (Qwen3-8B)         & 9.00B   & 10902.6 & 20.04\,GB & --           & 0.901 & 9.5   & 6.138 & 104.4 \\
    Student SWD-H (Qwen3-0.6B)   & 854.93M &   988.8 &  2.54\,GB & $11.0\times$ & 0.561 & 8.5   & 4.621 & 110.5 \\
     Student M-GRPO (Qwen3-0.6B)  & 854.93M &   988.8 &  2.54\,GB & $11.0\times$ & 0.523 & 7.9   & 3.105 & 70.8  \\
    \bottomrule
    \end{tabular}%
\end{table*}


For each training sample, we generate $G=4$ completions and compute group-relative advantages. This choice strikes a reasonable balance between group diversity and computational overhead.
\begin{equation}
A_i = r_i - \frac{1}{G}\sum_{j=1}^{G} r_j.
\end{equation}
We then optimize a PPO-style clipped objective with a KL penalty to the distilled reference model:
\begin{equation}
\mathcal{L}_{\mathrm{GRPO}}
=
-\mathbb{E}_i
\left[
\min\!\left(
\rho_i A_i,
\mathrm{clip}(\rho_i,1-\epsilon,1+\epsilon)A_i
\right)
\right]
+
\beta\,\mathrm{KL}(\pi_{\theta}\,\|\,\pi_{\mathrm{ref}}),
\end{equation}
where $\rho_i=\pi_{\theta}(y_i\mid x)/\pi_{\mathrm{old}}(y_i\mid x)$ and $\beta=0.1$. For stability, we clamp the log-ratio to $[-5,5]$ and use the non-negative KL form
\begin{equation}
\Delta_i=
\mathrm{clip}\!\left(
\log \pi_{\theta}(y_i\mid x)-\log \pi_{\mathrm{ref}}(y_i\mid x), -5, 5
\right),
\end{equation}
\begin{equation}
\mathrm{KL}_i=\exp(-\Delta_i)+\Delta_i-1.
\end{equation}
This stage encourages concise, complete, and emotion-faithful outputs, making the distilled student more efficient at inference time.


\section{Experimental Setup}

\subsection{Datasets}

We evaluate on nine benchmarks spanning three task families: basic emotion recognition (MER2023 \cite{MER2023}, MER2024 \cite{MER2024}, MELD \cite{MELD}, and IEMOCAP \cite{IEMOCAP}), sentiment analysis (CMU-MOSI \cite{CMU-MOSI}, CMU-MOSEI \cite{CMU-MOSEI}, CH-SIMS \cite{CH-SIMS}, and CH-SIMS v2 \cite{CH-SIMv2}), and fine-grained open-vocabulary MER (OV-MERD+ \cite{OV-MERD+}). This coverage is intentionally broad: a compact deployment model is only useful if it remains robust across fixed-label and open-set settings, sentiment and emotion tasks, and both Chinese and English data.

\subsection{Implementation Details}

The teacher uses Qwen3-8B, CLIP-ViT-Large-Patch14, and HuBERT-Large, and is trained for 60 epochs on MER-Caption+. The student uses Qwen3-0.6B, CLIP-ViT-Base-Patch16, and HuBERT-Base. All details are shown in supplementary materials.


\subsection{Evaluation Protocol}

For basic emotion recognition, we use HIT rate as the primary metric: a prediction is counted as correct if the ground-truth basic label appears in the predicted label set after emotion-wheel grouping. We additionally report \texttt{mscore} for these datasets, and defer their detailed analysis to the supplementary material. For sentiment analysis, we report weighted average F-score (WAF) and accuracy (ACC), using WAF as the primary metric because of label imbalance; detailed ACC analysis is also provided in the supplementary material. For fine-grained open-vocabulary MER, we follow the set-level formulation and report $\mathrm{Precision}_s$, $\mathrm{Recall}_s$, and $F_s$, averaged over five emotion wheels; $F_s$ is the primary metric. The dataset-wise mean reported in the main table averages the primary metric of each benchmark. 


\section{Experiments}

We organize the evaluation around five Research Questions.

\noindent \textbf{RQ1 (Effectiveness):} Can a sub-1B model match or surpass an 8B teacher across diverse MER benchmarks?

\noindent \textbf{RQ2 (Component Ablation):} What does each component (SWD-H, M-GRPO) contribute to the final performance?

\noindent \textbf{RQ3 (Distillation Design):} Among OT-based distillation variants, does SWD-H perform highest?

\noindent \textbf{RQ4 (Generation Refinement):} Does M-GRPO improve generation quality while maintaining recognition accuracy?

\noindent \textbf{RQ5 (Efficiency):} Does the compact model achieve practical deployment efficiency?

\begin{table*}[t]
    \centering
    \small
    \setlength{\tabcolsep}{4pt}
    \renewcommand{\arraystretch}{0.9}
    \caption{Module-removal ablation. The full model corresponds to the SWD distilled student with all reported modules enabled.}
    \vspace{-1em}
    \label{tab:module_ablation}
    \resizebox{\textwidth}{!}{
    \begin{tabular}{l|cccc|cccc|c| c}
        \toprule
        & \multicolumn{4}{c|}{Basic Emotion} & \multicolumn{4}{c|}{Sentiment} & \multicolumn{1}{c|}{Fine-grained} &  \\
        Variant & MER2023 & MER2024 & MELD & IEMOCAP & MOSI & MOSEI & SIMS & SIMS v2 & OV-MERD+ & Mean \\
        \midrule
        Student + SWD-H + M-GRPO (Ours)           & \underline{75.66} & \underline{80.04} & \underline{57.02} & \underline{61.46} & \textbf{80.04} & \textbf{77.23} & \textbf{89.05} & \textbf{87.74} & \textbf{63.28} & \textbf{74.61} \\
        w/o M-GRPO & \textbf{77.40} & \textbf{80.69} & \textbf{58.13} & \textbf{65.60} & \underline{76.14} & \underline{76.75} &\underline{ 86.52} & \underline{84.66} & \underline{61.59} & \underline{74.16} \\
        w/o SWD-H  & 70.67 & 76.78 & 54.66 & 59.87 & 72.48 & 73.33 & 85.00 & 84.29 & 58.37 & 70.61 \\
        \bottomrule
    \end{tabular}}
    \vspace{-1em}
\end{table*}

% -----------------------------
% Table 4
% -----------------------------
\begin{table*}[t!]
    \centering
    \small
    \setlength{\tabcolsep}{4pt}
    \renewcommand{\arraystretch}{0.9}
    \caption{Comparison of distillation variants. Among the completed runs, SWD achieves the best overall mean score. }
    \vspace{-1em}
    \label{tab:distill_ablation}
    \resizebox{\textwidth}{!}{
    \begin{tabular}{lccccccccc c}
        \toprule
        Method & MER2023 & MER2024 & MELD & IEMOCAP & MOSI & MOSEI & SIMS & SIMS v2 & OV-MERD+ & Mean \\
        \midrule
        CE only & 70.67 & 76.78 & 54.66 & 59.87 & 72.48 & 73.33 & 85.00 & 84.29 & 58.37 & 70.61 \\
        KL on logits only \cite{KL2014}& 75.82 & 78.67 & \underline{57.55} & \underline{65.32} & 75.60 & 76.57 & 85.91 & 84.65 & 61.28 & 73.49 \\
        \midrule
        Sinkhorn Optimal Transport (Logits) \cite{NIPS2013_af21d0c9} & 74.33 & 80.03 & 57.18 & 61.94 & 75.72 & 76.78 & 86.14 & 84.62 & 59.68 & 72.94 \\
        Sinkhorn Optimal Transport (Hidden) \cite{NIPS2013_af21d0c9} & \underline{76.44} & 79.57 & 56.80 & 62.58 & \textbf{76.99} & \underline{76.79} & 86.41 & \textbf{85.45} & \textbf{61.98} & \underline{73.67} \\
        Position-aware Optimal Transport \cite{PAOT} & 75.46 & 78.91 & 55.60 & 61.12 & 76.04 & 75.19 & 85.95 & 84.75 & 60.61 & 72.63 \\
        Unbalanced Optimal Transport \cite{UnbalancedOT} & 75.58 & \underline{80.05} & 57.19 & 64.53 & \underline{76.21} & 76.82 & \textbf{87.26} & 83.99 & 61.03 & 73.63 \\
        SWD (Logits) \cite{kolouri2019} & 74.26 & 78.86 & 56.22 & 62.40 & 75.76 & \textbf{78.91} & \underline{86.75} & \underline{84.76} & 61.27 & 73.24\\
        SWD-H (Ours)& \textbf{77.40} & \textbf{80.69} & \textbf{58.13} & \textbf{65.60} & 76.14 & 76.75 & 86.52 & 84.66 & \underline{61.59} & \textbf{74.16}\\
        \bottomrule
    \end{tabular}}
    \vspace{-1em}
\end{table*}

\subsection{Main Results (RQ1)}

In the full audio+video+text setting, Light-MER achieves a mean of 74.61, surpassing the 8B teacher at 73.93 by 0.68 points and the original AffectGPT at 69.77 by 4.84 points, all with a sub-1B student of 854M total parameters (Table~\ref{tab:main_results}). Our 8B teacher upgrades the AffectGPT backbone from Qwen2.5-7B to Qwen3-8B; the original 7B AffectGPT is listed separately. Light-MER leads on 7 of 9 benchmarks, spanning basic emotion, sentiment, and fine-grained recognition.

Light-MER also generalizes across modality subsets. Audio+text mean reaches 66.40 vs.\ teacher 65.07, and video+text mean reaches 65.71 vs.\ teacher 64.44. Across all three settings, the sub-1B student surpasses the 8B teacher on mean score while using $11\times$ fewer parameters (Table~\ref{tab:inference-profile}).

The student falls short of the teacher on MER2023 and CMU-MOSEI, both by roughly one point. On the remaining seven benchmarks, the distillation pipeline transfers multimodal patterns that compensate for the capacity gap (Section~\ref{sec:ablation}).

\textbf{(Answer to RQ1)}: A sub-1B student can match and surpass an 8B teacher in mean performance across diverse MER benchmarks when the distillation objective preserves the teacher's multimodal representation structure. The supplementary material further provides visualizations demonstrating the effectiveness of the distilled light-MER understanding.

\subsection{Ablation Study (RQ2)}
\label{sec:ablation}

Progressive addition on Table~\ref{tab:module_ablation} shows the contribution of each component. The CE-only student obtains a mean of 70.61, already 95.5\% of the teacher's 73.93. Adding SWD-H raises the mean to 74.16, a 3.55-point gain that accounts for 88.8\% of the total improvement to the full model at 74.61. Adding M-GRPO contributes a further 0.45 points.

SWD-H concentrates its gains on basic emotion tasks, with MER2023 improving by 6.73 points and IEMOCAP by 5.73, while sentiment benchmarks show smaller gains.  SWD-H aligns hidden-state distributions (Figure~\ref{fig:hidden_logit_analysis}d vs.\ panel~b) with the teacher and provides the primary accuracy gain. 
Furthermore, M-GRPO refines the student's generation behavior, producing concise and interpretable outputs suitable for deployment  (Section~\ref{sec:efficiency}). A slight drop in HIT on the basic emotion benchmarks is expected, since our reward design encourages shorter, more explanation-focused emotion descriptions to improve inference efficiency. Compared with the original verbose outputs, such concise descriptions are less likely to explicitly cover benchmark label words, which can slightly reduce HIT. The clear gains on WAF and $F_s$ confirm the overall effectiveness of the proposed design.

\textbf{(Answer to RQ2):} Both components are necessary. SWD-H provides the dominant accuracy gain through hidden-state alignment. M-GRPO complements SWD-H by refining generation quality and improving inference efficiency, producing shorter and more interpretable emotion descriptions for deployment.

\begin{figure}[t!]
    \centering
    % \vspace{-1.3em}
    \includegraphics[width=0.42\textwidth]{Image/loss_convergence_detailed1.pdf}
    \vspace{-1.3em}
    \caption{Loss convergence and decomposition for SWD-H distillation. (a) Under a fair comparison, the SWD-H distilled student shows faster, smoother cross-entropy convergence, reaching the target loss about 1.6$\times$ earlier than the teacher reference. The zoomed late-stage view highlights its stable optimization. (b) After OT weight ramp-up, the SWD-H term provides consistent geometric regularization while the total loss continues to decrease.}
    \label{fig:loss_convergence_detailed1}
    \vspace{-1.8em}
\end{figure}

\begin{figure}[t!]
    \centering
     \includegraphics[width=\linewidth]{Image/grpo_reward_phases.pdf}
     \vspace{-2.5em}
    \caption{Training dynamics of multi-reward GRPO optimization across four phases: exploration, rapid rise, fluctuation, and further refinement. The smoothed total reward increases steadily and stabilizes after approximately 500 steps, while the smoothed analysis length gradually decreases and converges to around 53 words by step 3200.}
    \label{fig:grpo_reward_phases}
    \vspace{-1.2em}
\end{figure}

\begin{table}[t!]
\centering
\renewcommand{\arraystretch}{0.9}
\caption{Description quality evaluated by GPT-5.4. We compare the student model before and after M-GRPO refinement across six dimensions.}
\vspace{-1em}
\label{tab:grpo-description-quality}
\begin{tabular}{lcccccc}
\toprule
\multirow{2}{*}{Stage} & \multicolumn{6}{c}{GPT-5.4~\cite{openai2026gpt54api}} \\
\cmidrule(lr){2-7}
 & Cls & Det & Flu & Sem & Con & Win\% \\
\midrule
Before M-GRPO   & 4.47 & 4.45 & 3.72 & 3.98 & 2.87 & 46 \\
After M-GRPO    & 4.89 & 3.62 & 4.24 & 3.83 & 4.07 & 54 \\
\bottomrule
\end{tabular}
\vspace{-1.2em}
\end{table}

\subsection{Distillation Variant Comparison (RQ3)}

Among eight distillation strategies (Table~\ref{tab:distill_ablation}), SWD-H achieves the highest overall mean of 74.16. All hidden-state OT methods except position-aware OT surpass KL-on-logits, which yields 73.49.

Two controlled comparisons isolate the source of improvement. Holding the loss function constant and switching from logits to hidden states improves Sinkhorn OT by 0.73 points, because hidden states encode how multimodal evidence is organized before decoding and carry richer distributional information than peaked output distributions (Section~\ref{subsec:observation}). Holding the alignment target at hidden states and switching from Sinkhorn to SWD adds another 0.49 points. SWD-H aligns teacher and student hidden states at answer-token positions. Each token contributes one hidden vector, so an answer of $L$ tokens gives only $L$ points to align. In MER, $L$ is typically 50--100, far fewer than the hidden dimension ($d=1024$). SWD projects these points onto random 1D directions and computes exact Wasserstein distances by sorting \cite{kolouri2019}, which works reliably even when $L$ is small. Sinkhorn OT \cite{NIPS2013_af21d0c9} solves an entropy-regularized transport plan iteratively, and the regularization strength must be set relative to $L$. With only 50--100 points, the plan becomes sensitive to this choice and the iterative updates can accumulate error across training steps. The convergence curves in Figure~\ref{fig:loss_convergence_detailed1} are consistent with this analysis, showing smoother convergence for SWD and oscillations for Sinkhorn.

The remaining OT variants each restrict how tokens are matched. Position-aware OT pairs each student token with the teacher token at the same sequence position, regardless of content, and yields the lowest mean among all OT variants. Unbalanced OT allows teacher tokens that have no similar student token to be left unmatched, and performs comparably to Sinkhorn. SWD on logits confirms the same pattern seen with Sinkhorn, with a similar margin of improvement when switching from logits to hidden states.

\textbf{(Answer to RQ3):} Two controlled comparisons identify the sources of SWD-H's advantage. Hidden-state alignment captures richer multimodal structure than logit-level matching. At the hidden-state level, SWD's closed-form sorting provides stable optimization on the short answer sequences typical of MER.


\subsection{M-GRPO Analysis (RQ4)}
\label{sec:mgrpo}

Adding M-GRPO refinement to the SWD-H-distilled student raises the mean from 74.16 to 74.61. As noted in Section~\ref{sec:ablation}, the gain concentrates on sentiment benchmarks while emotion benchmarks show trade-offs. The conciseness reward shortens the reasoning trace, which leaves fewer tokens for enumerating fine-grained emotion cues. Multi-class emotion recognition relies more on such enumeration than binary sentiment polarity, consistent with the observed accuracy shift.

GPT-5.4 evaluation of description quality (Table~\ref{tab:grpo-description-quality}) shows improvements on classification accuracy (Cls), fluency (Flu), and conciseness (Con), all scored on a five-point Likert-type scale \cite{Likert}, with conciseness showing the largest single gain (+1.20). Detail coverage (Det), scored on the same scale, drops, consistent with the shorter output, while semantic (Sem) coherence remains stable, suggesting that the core emotion reasoning is preserved as surface detail is compressed. In head-to-head comparisons, the M-GRPO-refined model wins 54\% of samples. Detailed prompt is provided in the supplementary material.

Optimizing five rewards simultaneously risks one reward dominating at the expense of others. The training dynamics in Figure~\ref{fig:grpo_reward_phases} show that both total reward and analysis length stabilize after roughly 500 steps without such collapse, with analysis length settling near 53 words by step 3200. 

\textbf{(Answer to RQ4):} M-GRPO improves aggregate recognition accuracy (+0.45 mean), generation fluency, and conciseness at the cost of reduced detail coverage. For deployment, shorter and focused descriptions are preferable to exhaustive multimodal analyses.

\subsection{Inference Efficiency (RQ5)}
\label{sec:efficiency}

Light-MER supports two inference modes. In \emph{direct} mode, the model outputs label tokens without intermediate reasoning. In \emph{descriptive} mode, the model generates an emotion description, from which a label extractor derives the final prediction. Table~\ref{tab:inference-profile} profiles model size, memory footprint, and per-sample latency for both modes.

The sub-1B student achieves $11.0\times$ compression in parameters and FLOPs, and $7.9\times$ memory reduction from 20.04\,GB to 2.54\,GB. In descriptive mode, M-GRPO refinement reduces the average output from 110.5 to 70.8 words and latency from 4.6\,s to 3.1\,s, a 32.8\% reduction over the SWD-H-only student. The latency gain comes entirely from learned brevity (Section~\ref{sec:mgrpo}), not architectural modification. In direct mode, the M-GRPO student runs at 0.5\,s per sample. Direct mode eliminates the reasoning trace, reducing latency at the cost of the explicit evidence chain needed for open-set label extraction. For closed-set tasks where the label space is fixed, direct mode offers a viable low-latency alternative. 

\textbf{(Answer to RQ5):} The Light-MER student fits within 2.54\,GB memory and processes each sample in 0.5\,s in direct mode or 3.1\,s in descriptive mode, with $11\times$ fewer parameters than the teacher. The efficiency profile shows that competitive generative MER is achievable with sub-1B models under the proposed distillation-refinement pipeline.


\section{Conclusion}
In this study, we revisited generative multimodal emotion recognition from an efficiency perspective and asked whether high-quality MER truly requires deployment models larger than 1B parameters. We have shown  that it is not necessary. We presented Light-MER, a lightweight teacher-student framework that transfers multimodal emotion reasoning from a large MLLM to a sub-1B deployment model through geometry-aware hidden-state distillation and post-distillation policy refinement. Specifically, SWD-H preserves the latent structure of multimodal representations during compression, while M-GRPO further improves generation quality and efficiency with multi-reward optimization. Extensive experiments on nine benchmarks show that Light-MER not only matches but surpasses an 8B teacher in mean performance, while significantly reducing memory usage and inference latency. These results suggest that efficient generative MER does not require large deployment models, and that preserving representation geometry is a promising direction for compact multimodal reasoning systems.

\balance
\bibliographystyle{ACM-Reference-Format}
\bibliography{sample-base}

\appendix

\end{document}



